\section{Reliable delivery}
\label{sec:reliability}

The checksum gives \textbf{integrity}: whoever receives a frame knows that it is intact. It gives
no \textbf{delivery guarantee}: the sender still does not know whether anything arrived.

That difference is this whole section, and the four mechanisms that follow build on each other in
turn.

\subsection{ACK and NACK}

The receiver answers. If the frame arrived intact it sends an \code{Ack}; if it was broken in a
way that can still be answered, it sends a \code{Nack}.

The rules for the acknowledgement frame:

\begin{itemize}
  \item \code{DST} and \code{SRC} swap places compared with the acknowledged frame.
  \item \code{SEQ} is the \textbf{same} as in the acknowledged frame --- that is how the sender
    knows what is being acknowledged.
  \item \code{LEN} is zero.
\end{itemize}

That \code{SEQ} comes along is not a detail. Without it, a sender with two frames in flight does
not know which of them arrived.

\begin{aside}
\note A \code{Nack} can only be sent when the frame is broken in a way that leaves the sender and
the sequence number readable. If the checksum is wrong, \emph{no} field can be trusted ---
including \code{SRC} --- so in practice a missing acknowledgement is the usual way an error shows
itself.
\end{aside}

\subsection{Timeout}

And with that: what does the sender do when nothing comes back?

It cannot wait forever, and it cannot tell the difference between a lost frame, a lost
acknowledgement and a slow receiver. The only way out is to make up its mind after a certain
time: a \textbf{timeout}.

Choosing the timeout value is a trade-off with no right answer:

\begin{narrowtable}{@{}ll@{}}
\thead{Too short} & \thead{Too long} \\ \midrule
Unnecessary retransmissions & Slow error detection \\
More duplicates & The bus stands still and waits \\
The bus is loaded more when it is already loaded & The system feels hung \\
\end{narrowtable}

\subsection{Retransmission}

On a timeout the frame is sent again, with the \textbf{same} \code{SEQ}. That is crucial: a new
sequence number would make the retransmission a new message, and the receiver would not be able
to recognise it as a repeat.

The number of attempts has to be limited. A sender that tries for all eternity against a node
that is switched off does nothing else.

\subsection{Duplicates}

And here the circle closes. Retransmission is exactly what creates duplicates.

The dangerous case is not that the frame is lost --- then nothing happened --- but that the
\textbf{acknowledgement} is lost:

\begin{console}
A -> B   StatusRequest, SEQ 0x0007      B receives and processes it
B -> A   Ack, SEQ 0x0007                the acknowledgement is lost
A        timeout
A -> B   StatusRequest, SEQ 0x0007      retransmission
B                                       ...processes it a second time?
\end{console}

From B's point of view these are two identical frames. If it processes both, the command has been
carried out twice, and if the command is "increase the setpoint by one" the result is wrong.

\textbf{The solution is to remember.} Every node stores the most recently processed \code{SEQ}
per sender. If the same \code{SEQ} turns up again:

\begin{itemize}
  \item Do \textbf{not} run the application logic again.
  \item Send the acknowledgement again \textbf{anyway}.
\end{itemize}

The second point is the one that gets forgotten. The sender is still waiting for its
acknowledgement, and silence only leads to yet another retransmission.

\begin{aside}
\note Four mechanisms, in a chain where each solves the problem the previous one created: the
checksum detects errors, the ACK confirms delivery, the timeout detects a missing ACK,
retransmission fixes it --- and duplicate detection cleans up after the retransmission. None of
them can be left out, and that is why the chain is worth knowing by heart.
\end{aside}

\subsection{And then the hardware does it for us}

Everything in this section is code you write, test and maintain. From \kapref{ch:can} onwards
most of it sits in silicon: CAN does framing, checksumming and collision handling in hardware,
and acknowledges in a single bit inside the frame as well.

That does not make this chapter less important. It makes it the background against which CAN can
be understood --- and a part of what the written test examines.
